\documentclass[aspectratio=169,11pt]{beamer}
\usepackage{fontspec}
\usepackage{listings}
\usepackage{xcolor}
\usepackage{booktabs}
\usepackage{array}

% ================= Palette (identique au livre Tome 1) =================
\definecolor{PyBleu}{HTML}{1A4A7A}
\definecolor{PyBleuMoyen}{HTML}{2D7AB3}
\definecolor{PyBleuClair}{HTML}{EAF1F8}
\definecolor{PyJaune}{HTML}{FFD43B}
\definecolor{PyNuit}{HTML}{142840}
\definecolor{PyTexte}{HTML}{1F2A37}
\definecolor{PyGris}{HTML}{5A6B7E}
\definecolor{PyRouge}{HTML}{B3442D}
\definecolor{PyVert}{HTML}{2E7D5B}
\definecolor{CodeTexte}{HTML}{E8EEF6}

% ================= Code Python =================
\lstdefinestyle{py}{
  language=Python,
  keywordstyle=\color{PyJaune}\bfseries,
  stringstyle=\color{PyJaune},
  commentstyle=\color{PyGris}\itshape,
  showstringspaces=false,
  breaklines=true,
  backgroundcolor=\color{PyNuit},
  basicstyle=\ttfamily\footnotesize\color{CodeTexte},
  columns=fullflexible,
  keepspaces=true,
  xleftmargin=6pt,
  framexleftmargin=6pt,
}
\lstdefinestyle{pysmall}{
  language=Python,
  keywordstyle=\color{PyJaune}\bfseries,
  stringstyle=\color{PyJaune},
  commentstyle=\color{PyGris}\itshape,
  showstringspaces=false,
  breaklines=true,
  backgroundcolor=\color{PyNuit},
  basicstyle=\ttfamily\scriptsize\color{CodeTexte},
  columns=fullflexible,
  keepspaces=true,
  xleftmargin=5pt,
  framexleftmargin=5pt,
}

% ================= Titres de diapositives =================
\setbeamercolor{structure}{fg=PyBleu}
\setbeamerfont{frametitle}{series=\bfseries, size=\Large}
\setbeamercolor{frametitle}{bg=PyBleu, fg=white}
\setbeamertemplate{frametitle}{%
  \nointerlineskip
  \begin{beamercolorbox}[wd=\textwidth, ht=2.1ex, dp=0.9ex,
                          leftskip=0.35cm, rightskip=0.35cm]{frametitle}%
    \usebeamerfont{frametitle}\insertframetitle%
  \end{beamercolorbox}%
  \vspace*{-0.4pt}%
  {\color{PyJaune}\rule{\textwidth}{1.6pt}}%
}

% ================= Pied de page =================
\setbeamerfont{footline}{size=\tiny}
\setbeamercolor{footline}{bg=white, fg=PyGris}
\newcommand{\PiedDePage}{Python en 24 heures \textbar\ Séance 09 : Héritage et méthodes spéciales}
\setbeamertemplate{footline}{%
  \leavevmode%
  \hbox{%
    \begin{beamercolorbox}[wd=\paperwidth, ht=2.2ex, dp=1.1ex,
                            leftskip=0.35cm, rightskip=0.35cm]{footline}%
      \usebeamerfont{footline}\PiedDePage
      \hfill\insertframenumber\,/\,\inserttotalframenumber%
    \end{beamercolorbox}}%
}

\setbeamertemplate{navigation symbols}{}
\setbeamertemplate{itemize item}{\color{PyBleuMoyen}$\blacktriangleright$}
\setbeamertemplate{itemize subitem}{\color{PyBleuMoyen}\textbullet}

% ================= Blocs pedagogiques =================
\setbeamercolor{block title}{bg=PyBleu, fg=white}
\setbeamercolor{block body}{bg=PyBleuClair, fg=PyTexte}
\setbeamercolor{block title alerted}{bg=PyRouge, fg=white}
\setbeamercolor{block body alerted}{bg=PyRouge!9, fg=PyTexte}
\setbeamercolor{block title example}{bg=PyJaune, fg=PyNuit}
\setbeamercolor{block body example}{bg=PyJaune!14, fg=PyTexte}

% ================= Metadonnees =================
\hypersetup{
  pdftitle={Python en 24 heures : Séance 09 : Héritage et méthodes spéciales},
  pdfauthor={HOUNKPE Mawulikplimi Roland et DJABON Ounimborbitibou},
  pdfsubject={De l'algorithmique à la programmation orientée objet},
  pdfkeywords={Python, algorithmique, POO, formation, exercices corrigés}
}

\begin{document}

{%
\setbeamercolor{background canvas}{bg=PyNuit}
\begin{frame}[plain]
\vspace{10mm}
{\color{PyJaune}\small\bfseries FORMATION « PYTHON EN 24 HEURES » \quad\textbar\quad SÉANCE 09\par}
\vspace{4mm}
{\color{white}\Huge\bfseries Héritage, polymorphisme et méthodes spéciales\par}
\vspace{4mm}
{\color{PyBleuClair}\large Chapitres 17 et 18 du Tome 1\par}
\vspace{7mm}
{\color{PyJaune}\rule{34mm}{2.2pt}\par}
\vspace{7mm}
{\color{white}\normalsize 1 h de cours (présentation Beamer) \quad$\bullet$\quad 1 h de pratique (notebook Jupyter)\par}
\vspace{2.5mm}
{\color{PyBleuClair}\normalsize HOUNKPE Mawulikplimi Roland et DJABON Ounimborbitibou\par}
\end{frame}
}
\begin{frame}{Objectifs de la séance}
\begin{itemize}
  \item Faire hériter une classe d'une autre et étendre \_\_init\_\_ avec super()
  \item Redéfinir des méthodes et exploiter le polymorphisme
  \item Choisir entre héritage (EST UN) et composition (A UN)
  \item Implémenter \_\_repr\_\_, \_\_eq\_\_, \_\_add\_\_, \_\_len\_\_ et \_\_getitem\_\_
  \item Comprendre l'ordre MRO en cas d'héritage multiple
\end{itemize}
\end{frame}
\begin{frame}{Rappel de la séance 08}
Vos classes savent naître, se protéger et se présenter : elles vont maintenant se transmettre des comportements et dialoguer avec les opérateurs de Python.

\begin{itemize}
  \item \_\_init\_\_ et self : la naissance et l'identité de l'instance
  \item Attributs d'instance contre attributs de classe partagés
  \item Encapsulation : \_interne, \_\_privé, @property avec validation
  \item @classmethod pour les fabriques, @staticmethod pour les utilitaires
  \item \_\_str\_\_ : la première méthode spéciale rencontrée
\end{itemize}
\end{frame}
\begin{frame}[fragile]{L'héritage : réutiliser sans recopier}
\begin{columns}[t]
\begin{column}{0.42\textwidth}
\begin{itemize}
  \setlength{\itemsep}{1.5pt}\setlength{\parskip}{0pt}
  \item class Voiture(Vehicule) : Voiture \textbf{hérite} de tout Vehicule
  \item Voiture est un Vehicule spécialisé : relation \textbf{EST UN}
  \item Les méthodes héritées fonctionnent sans être réécrites
  \item On peut \textbf{redéfinir} une méthode pour la spécialiser
  \item super() appelle la version de la classe parente
\end{itemize}
\end{column}
\hspace{0.02\textwidth}
\begin{column}{0.55\textwidth}
\begin{lstlisting}[style=pysmall]
class Vehicule:
    def __init__(self, marque):
        self.marque = marque
    def klaxon(self):
        return "Tut tut"

class Voiture(Vehicule):
    def klaxon(self):
        return "Pouet pouet"

v = Voiture("Toyota")
print(v.marque, v.klaxon())
\end{lstlisting}
\end{column}
\end{columns}
\end{frame}

\begin{frame}[fragile]{super().\_\_init\_\_() : étendre sans casser}

\begin{lstlisting}[style=py]
class Vehicule:
    def __init__(self, marque, roues):
        self.marque = marque
        self.roues = roues

class Moto(Vehicule):
    def __init__(self, marque, cylindree):
        super().__init__(marque, roues=2)
        self.cylindree = cylindree

m = Moto("Yamaha", 600)
print(m.marque, m.roues, m.cylindree)
\end{lstlisting}
\vspace{0.5pt}

{\small\color{PyGris}\itshape super() délègue au parent la partie commune, puis Moto ajoute la sienne : zéro duplication du code d'initialisation.}

\end{frame}

\begin{frame}{Polymorphisme : même appel, comportements variés}
\begin{itemize}
  \item Une boucle peut appeler la même méthode sur des objets de classes différentes
  \item Chaque classe répond à sa façon : c'est le \textbf{polymorphisme}
  \item isinstance(x, Classe) teste la filiation en respectant l'héritage
  \item Le code appelant ne connaît que le contrat : `forme.aire()`
  \item Ajouter une nouvelle forme ne modifie \textbf{aucune} ligne du code appelant
\end{itemize}
\end{frame}

\begin{frame}[fragile]{isinstance : tester la filiation}
\begin{columns}[t]
\begin{column}{0.42\textwidth}
\begin{itemize}
  \setlength{\itemsep}{1.5pt}\setlength{\parskip}{0pt}
  \item `isinstance(x, Vehicule)` vaut True pour x Vehicule OU x Voiture
  \item L'héritage traverse le test : une fille EST UN parent
  \item Attention : trop de isinstance, c'est du code procédural déguisé
  \item Préférez le polymorphisme : appelez la méthode, la classe répond
  \item isinstance reste légitime aux frontières : validation, sérialisation
\end{itemize}
\end{column}
\hspace{0.02\textwidth}
\begin{column}{0.55\textwidth}
\begin{lstlisting}[style=pysmall]
formes = [Cercle(1), Rectangle(2, 3)]

for forme in formes:
    print(isinstance(forme, Forme))

# le polymorphisme remplace les tests :
total = sum(f.aire() for f in formes)
\end{lstlisting}
\end{column}
\end{columns}
\end{frame}

\begin{frame}[fragile]{Composition : la relation A UN}

\begin{lstlisting}[style=py]
class Moteur:
    def __init__(self, puissance):
        self.puissance = puissance

class Voiture:                    # sans héritage !
    def __init__(self, marque, puissance):
        self.marque = marque
        self.moteur = Moteur(puissance)   # A UN moteur

v = Voiture("Toyota", 120)
print(v.moteur.puissance)
\end{lstlisting}
\vspace{0.5pt}

{\small\color{PyGris}\itshape Une Voiture n'est pas un Moteur : elle en possède un. Héritage pour EST UN, composition pour A UN : la question à se poser avant chaque classe.}

\end{frame}

\begin{frame}[fragile]{\_\_repr\_\_ et \_\_str\_\_ : deux vitrines}

\begin{lstlisting}[style=py]
class Duree:
    def __init__(self, heures, minutes):
        self.total = heures * 60 + minutes

    def __str__(self):
        return f"{self.total // 60} h {self.total % 60:02d}"

    def __repr__(self):
        return f"Duree({self.total // 60}, {self.total % 60})"

print(Duree(1, 90))        # 2 h 30
print([Duree(1, 90)])      # [Duree(2, 30)]
\end{lstlisting}
\vspace{0.5pt}

{\small\color{PyGris}\itshape \_\_str\_\_ pour l'utilisateur final (print), \_\_repr\_\_ pour le développeur (console, listes, débogueur). Un \_\_repr\_\_ honnête montre comment recréer l'objet.}

\end{frame}

\begin{frame}[fragile]{\_\_eq\_\_ et \_\_lt\_\_ : comparer}

\begin{lstlisting}[style=py]
class Duree:
    def __init__(self, heures, minutes=0):
        self.total = heures * 60 + minutes

    def __eq__(self, autre):
        return self.total == autre.total

    def __lt__(self, autre):
        return self.total < autre.total

print(Duree(1, 90) == Duree(2, 30))   # True
print(sorted([Duree(2, 0), Duree(1, 0)]))
\end{lstlisting}
\vspace{0.5pt}

{\small\color{PyGris}\itshape Définir \_\_eq\_\_ rend == correct ; définir \_\_lt\_\_ suffit à sorted() pour trier les objets.}

\end{frame}

\begin{frame}[fragile]{\_\_len\_\_ et \_\_getitem\_\_ : se comporter comme une liste}

\begin{lstlisting}[style=py]
class Playlist:
    def __init__(self, titres):
        self.titres = list(titres)

    def __len__(self):
        return len(self.titres)

    def __getitem__(self, index):
        return self.titres[index]

p = Playlist(["Awa", "Kofi"])
print(len(p))        # 2
print(p[0], p[-1])   # Awa Kofi
for titre in p:      # le for marche aussi !
    print(titre)
\end{lstlisting}
\vspace{0.5pt}

{\small\color{PyGris}\itshape len() appelle \_\_len\_\_, les crochets appellent \_\_getitem\_\_ : une classe qui implémente ces deux méthodes devient iterable et indexable gratuitement.}

\end{frame}

\begin{frame}[fragile]{Héritage multiple et MRO}
\begin{columns}[t]
\begin{column}{0.42\textwidth}
\begin{itemize}
  \setlength{\itemsep}{1.5pt}\setlength{\parskip}{0pt}
  \item Une classe peut hériter de plusieurs parents : class C(A, B)
  \item Python résout un nom en remontant les parents dans un ordre précis
  \item Cet ordre est le \textbf{MRO} : Method Resolution Order
  \item C.\_\_mro\_\_ révèle la chaîne de recherche
  \item Les mixins offrent des comportements optionnels par héritage multiple léger
\end{itemize}
\end{column}
\hspace{0.02\textwidth}
\begin{column}{0.55\textwidth}
\begin{lstlisting}[style=pysmall]
class Rechargeable:
    def recharger(self):
        return "batterie pleine"

class Roule:
    def rouler(self):
        return "je roule"

class Trottinette(Rechargeable, Roule):
    pass

t = Trottinette()
print(t.recharger(), t.rouler())
print(Trottinette.__mro__)
\end{lstlisting}
\end{column}
\end{columns}
\end{frame}

\begin{frame}{Attention}
\begin{alertblock}{Le MRO et le diamant}
En héritage multiple, deux parents peuvent définir la même méthode : c'est le premier de la liste \_\_mro\_\_ qui gagne, pas le plus « logique ». Évitez les hiérarchies en diamant tant que le MRO n'est pas limpide pour vous, et préférez la composition quand l'héritage n'exprime pas clairement EST UN.
\end{alertblock}
\end{frame}

\begin{frame}{À retenir}
\begin{block}{L'essentiel de la séance 09}
L'héritage exprime EST UN et se complète par super() ; la composition exprime A UN et reste souvent plus souple. Le polymorphisme fait parler chaque classe avec sa voix. Les méthodes spéciales (\_\_repr\_\_, \_\_eq\_\_, \_\_lt\_\_, \_\_len\_\_, \_\_getitem\_\_) branchent vos objets sur le langage lui même : print, ==, sorted, len et les crochets.
\end{block}
\end{frame}

\begin{frame}{Travaux pratiques (1 h) : Héritage et méthodes spéciales}
Cinq exercices pour brancher vos classes sur le langage. L'exercice 5 construit une petite bibliothèque de médias complète : gardez la, elle inspire le mini projet de la séance 10.

\begin{itemize}
  \setlength{\itemsep}{1.5pt}\setlength{\parskip}{0pt}
  \item \textbf{Exercice 1 (Basique)} : Vehicule et Voiture
  \item \textbf{Exercice 2 (Basique)} : \_\_str\_\_ et \_\_repr\_\_ sur Media
  \item \textbf{Exercice 3 (Intermédiaire)} : Formes polymorphes
  \item \textbf{Exercice 4 (Intermédiaire)} : Classe Duree avec \_\_add\_\_ et \_\_eq\_\_
  \item \textbf{Exercice 5 (Défi)} : Bibliothèque de médias indexable
\end{itemize}
\vspace{4pt}

{\small\color{PyGris}Ouvrez le notebook « Seance\_09\_heritage\_et\_methodes\_speciales.ipynb » pour commencer.}
\end{frame}
\begin{frame}{Pour la prochaine séance}
\begin{exampleblock}{Conseil de travail}
Séance 10 : le mini projet bancaire. Nous allons assembler trois classes qui se parlent, avec historique et validations. Relisez CompteSimple de la séance 08 : c'est son jumeau évolué qui ouvre le TP.
\end{exampleblock}
\end{frame}

\end{document}
